\chapter{Por qué \emph{khop\_reach} delta es difícil}
\label{cap:khop-delta-hard}

\begin{intuicion}
Stage 2a (capítulo~\ref{cap:incremental-nlf}) resolvió el caso
incremental para NLF: cada \texttt{add\_relation} bumpea un
contador local. Para k-hop reachability el problema es
fundamentalmente más difícil porque la información NO es local. Una
sola arista \((u, v)\) puede cambiar el bit de \(v\) ---trivial---
pero también el de \emph{todos los predecesores} de \(u\) que
acortan su camino al tipo de \(v\). En grafos scale-free con hubs,
``todos los predecesores'' puede ser millones de vértices. Este
capítulo argumenta por qué la solución no es invalidar-y-rebuild
ni una propagación naive, y prepara el terreno para Delta WCOJ
Propagation (próximo capítulo).
\end{intuicion}

\section{El problema en una frase}

NLF es \emph{local}: una arista \((u, v)\) afecta solamente
\(\NLF(u, \mathrm{type}(v))\). Update incremental: bumpear un
contador. \(O(1)\) por arista.

k-hop es \emph{transitivo}: una arista \((u, v, t)\) puede cambiar
\(\bits[w]\) y \(\mathrm{min/max\_first}[w][L]\) para todo \(w\)
que tiene un camino a \(u\) en \(< \texttt{MAX\_DEPTH}\) saltos.
``Todo \(w\)'' puede ser cero (si la arista no agrega tipos
nuevos), pocos, o muchísimos.

\section{Ejemplo: una arista, muchos cambios}

\begin{ejemplo}
Grafo previo a la inserción:
\begin{itemize}
    \item \(u_1, u_2, u_3 \to h\) (tres usuarios autenticando en un
    host).
    \item \(h \to ip_1\) en \(\tau = 100\) (host conectado a una IP).
\end{itemize}

\noindent
Para \texttt{MAX\_DEPTH = 2}, los \(u_i\) tienen \(\bits[u_i]\)
con bits \texttt{Host} (vía \(h\)) e \texttt{IP} (vía \(h \to
ip_1\)) prendidos. \(\mathrm{min/max\_first}[u_i][\texttt{IP}] =
\tau\) de la arista \(u_i \to h\).

Llega una arista nueva: \(h \to ip_2\) en \(\tau = 200\).

\textbf{Cambios al índice}:
\begin{itemize}
    \item \(\bits[h]\) ya tiene bit \texttt{IP} prendido (por \(ip_1\)).
    Sin cambio.
    \item \(\mathrm{max\_first}[h][\texttt{IP}]\) actualiza a 200
    (era 100).
    \item Para cada \(u_i\): \(\bits[u_i]\) ya tiene IP. Sin cambio
    en bits.
    \item \(\mathrm{max\_first}[u_i][\texttt{IP}]\) actualiza al
    timestamp de la arista \(u_i \to h\) si es posterior. (Depende
    del análisis de paths concretos.)
\end{itemize}

\noindent
Tres mutaciones de min/max\_first, ninguna en bits. Manageable.

Ahora considerá una arista que SÍ agrega tipo nuevo: digamos un
\(d\) (Domain) y la arista \(h \to d\) en \(\tau = 300\).

\textbf{Cambios}:
\begin{itemize}
    \item \(\bits[h]\) ahora tiene bit \texttt{Domain} (era 0).
    Cambio en bits.
    \item Por transitividad: para cada \(u_i\), \(\bits[u_i]\) ahora
    tiene bit \texttt{Domain} también. Tres cambios en bits.
    \item Si los \(u_i\) tienen predecesores propios (digamos cada
    uno tiene un grupo de \(\sim 1000\) predecesores), \emph{esos
    también} ganan el bit \texttt{Domain}. Cascada: \(1000 \cdot 3 =
    3000\) cambios adicionales.
    \item Si el grafo es scale-free y \(h\) es un hub, los
    predecesores y predecesores-de-predecesores pueden ser
    \(\sim 10^6\). Cascada explota.
\end{itemize}
\end{ejemplo}

\section{Por qué invalidar-y-rebuild es lento}

La política pre-stage-2a era simple: cualquier
\texttt{add\_relation} \emph{post}-build setea
\texttt{self.khop\_reach = OnceLock::new()}; el próximo hunt
reconstruye desde cero.

Costo: \(O(|\Lc_V| \cdot \texttt{MAX\_DEPTH} \cdot |\E|)\)
(capítulo~\ref{cap:khop-detalle}). Para \(|\E| = 10^8\): $\sim$1.5
minutos por rebuild. Para ingestas continuas (Sentinel pulling
cada 30 s + 1 hunt por minuto): rebuild cada 30 s. La UI vive en
estado de espera permanente.

\section{Por qué la propagación naive falla}

La idea natural: cuando llega \((u, v, t)\), recursar hacia atrás
desde \(u\) actualizando \(\bits\) y \(\mathrm{min/max\_first}\)
de cada predecesor:

\begin{lstlisting}
fn propagate_naive(u: StrId, delta_bits: u16, t: i64) {
    bits[u] |= delta_bits;
    update_min_max(u, t);
    for w in predecessors(u) {
        propagate_naive(w, delta_bits, t);  // recurse
    }
}
\end{lstlisting}

Tres problemas:

\begin{enumerate}
    \item \textbf{Sin cota de profundidad}. Recursión hasta agotar
    \emph{todos} los ancestros. Para grafo conexo, eso es \(|\V|\).
    Pero el índice sólo cuenta caminos de \(\leq \texttt{MAX\_DEPTH}\)
    saltos; propagación más allá es work pero no aporta info.
    \item \textbf{Sin dedup}. El mismo \(w\) puede ser predecesor
    por varios paths; la propagación naive lo visita cada vez. Costo
    explota.
    \item \textbf{Sin chequeo de delta vacío}. Si \(\bits[w]\) ya
    tenía todos los bits que \(\bits[u]\) está agregando, la
    propagación sigue para nada.
\end{enumerate}

\noindent
Cada problema requiere su antídoto: cota por
\texttt{MAX\_DEPTH}, frontier set para dedup, chequeo
\(\Delta = \bits[u] \setminus \bits[w]\) antes de propagar.

\section{Por qué min/max\_first es aún más difícil}

\(\bits\) es monotónico-creciente: nunca se quita un bit; sólo
se agregan. Una vez setado, queda. La sobre-aproximación está
intrínsecamente protegida.

\(\mathrm{min\_first}\) y \(\mathrm{max\_first}\) son
\emph{bidireccionales}: una arista nueva podría \emph{bajar}
\(\mathrm{min\_first}\) o \emph{subir} \(\mathrm{max\_first}\). Pero
si una arista vieja se \emph{borra} (caso que GraphHunter no
maneja hoy ---el grafo es append-only, capítulo
\ref{cap:entity-relation}), las cantidades extremales podrían
\emph{volverse menos extremas} ---requiriendo recomputar.

Por suerte, append-only nos salva: nunca borramos, así que
\(\mathrm{min\_first}\) sólo baja y \(\mathrm{max\_first}\) sólo
sube. Mantener esa invariante bajo deltas es manejable. Pero
\emph{detectar} cuándo cambian para un \(w\) dado requiere
re-mirar el path completo desde \(w\) ---no sólo el bit.

\section{Soundness es la salvaguarda}

Hay una buena noticia: \emph{cualquier} sobre-aproximación es
sound. El motor podría:

\begin{itemize}
    \item \textbf{No actualizar nunca}: el bit existente sigue
    válido (no se quita); el \(\mathrm{min/max\_first}\) sigue
    siendo extremal según lo que conocía. Nuevo info no se refleja
    pero las podas que aplicaba siguen siendo correctas. La
    ``pérdida'' es que algunos prunes que ahora podrían descartar
    quedan aceptando ---el matcher trabaja un poco más, no menos.
    \item \textbf{Actualizar parcialmente}: actualizar \(\bits\) y
    \(\mathrm{min/max\_first}\) sólo para \(u\) y \(v\)
    directamente. Predecesores quedan con info vieja
    (sobre-aproximada). El matcher sigue siendo correcto.
\end{itemize}

\noindent
La razón por la que el motor \emph{quiere} actualizar todo lo más
posible: cada bit \(L\) extra prendido en \(\bits[w]\) significa
que \(w\) sobrevive un prune que antes lo descartaba ---más
trabajo para el matcher---. La calidad del prune depende de qué
tan tight sea la sobre-aproximación.

\section{El plan: Delta WCOJ Propagation}

El próximo capítulo desarrolla la solución que el motor adoptará
(stage 2b futuro): \emph{Delta WCOJ Propagation}, basado en
Kankanamge-Sahu-Mhedhbi-Salihoglu SIGMOD 2017. La idea de fondo:

\begin{enumerate}
    \item Computar el delta \(\Delta_u = \bits[v] \setminus
    \bits[u]\). Si vacío, abortar (la arista no aporta info nueva
    desde \(u\)).
    \item Si no vacío, actualizar \(\bits[u]\) y propagar el
    \emph{delta} (no el bits completo) hacia atrás vía reverse
    adjacency.
    \item En cada nivel, aplicar de nuevo el chequeo
    ``delta vacío $\Rightarrow$ abortar esta rama''.
    \item Cota \texttt{MAX\_DEPTH} en la cascada.
    \item \textbf{Cascade budget}: si la cascada visita más de un
    threshold (default 100K vértices), abortar y caer a
    invalidación. Caso patológico de hubs.
\end{enumerate}

\noindent
Costo amortizado: \(O(\texttt{MAX\_DEPTH} \cdot
d_{\mathrm{reverse}})\) por arista, donde \(d_{\mathrm{reverse}}\)
es el grado promedio del grafo inverso. Para grafos scale-free
típicos, \(\sim 10\)--\(100\times\) más barato que rebuild.

\section{Resumen}

\begin{itemize}
    \item NLF es local; k-hop es transitivo. Esa diferencia hace
    que el incremental update de k-hop sea fundamentalmente más
    difícil.
    \item Una sola arista puede generar cascada de cambios
    proporcional al grado \emph{reverse} de los predecesores
    indirectos.
    \item Invalidate-and-rebuild es lento: $\sim$1.5 min para 28 GB
    EVTX.
    \item Propagación naive falla por falta de cota de profundidad,
    falta de dedup, y falta de chequeo delta vacío.
    \item \(\mathrm{min/max\_first}\) son más sutiles que
    \(\bits\) por bidireccionalidad, pero append-only nos salva.
    \item Soundness garantiza que cualquier sobre-aproximación
    funciona; ``actualizar más'' sólo mejora la calidad del
    prune.
    \item Delta WCOJ Propagation (próximo capítulo) ataca el
    problema con propagación reverse acotada por
    \texttt{MAX\_DEPTH} y \emph{cascade budget}.
\end{itemize}

\begin{ejercicios}
\begin{enumerate}
    \item \starred{} Para un grafo scale-free con hub de
    in-degree \(10^4\), si llega una arista nueva \((u, v)\)
    donde \(v\) tiene tipo no visto antes y \(u\) es predecesor
    inmediato del hub, ¿cuántos vértices podrían cambiar
    \(\bits\)? Estimación.

    \item ¿Por qué la propagación naive sin dedup explota? Pensá
    en un grafo bipartito: cada vértice de un lado tiene 100
    vecinos del otro, vice versa.

    \item \starred{} Demostrá: si una arista nueva NO aporta tipos
    nuevos (\(\bits[v] \subseteq \bits[u]\)), no hace falta
    propagar bits. ¿Y min/max\_first?

    \item Imaginá una variante donde \(\bits\) usa más bits (32 en
    vez de 16). ¿Cambia la facilidad del incremental? ¿Por qué?

    \item Si el motor permitiera DELETE de aristas, ¿qué se rompe
    en la invariante append-only? ¿Cómo lo arreglarías?
\end{enumerate}
\end{ejercicios}
